\documentclass[12pt,a4paper]{article}

% ─────────────────────────────────────────────
%  PACKAGES
% ─────────────────────────────────────────────
\usepackage[margin=1in]{geometry}
\usepackage{amsmath,amssymb,amsthm}
\usepackage{xcolor}
\usepackage{tikz}
\usetikzlibrary{arrows.meta,shapes,positioning,fit,backgrounds,decorations.pathreplacing,calc}
\usepackage{tcolorbox}
\tcbuselibrary{skins,breakable,theorems}
\usepackage{hyperref}
\usepackage{enumitem}
\usepackage{booktabs}
\usepackage{array}
\usepackage{graphicx}
\usepackage{fancyhdr}
\usepackage{titlesec}
\usepackage{mdframed}
\usepackage{algorithm}
\usepackage{algpseudocode}
\usepackage{multicol}
\usepackage{float}

% ─────────────────────────────────────────────
%  COLORS
% ─────────────────────────────────────────────
\definecolor{darkblue}{RGB}{0,51,102}
\definecolor{midblue}{RGB}{0,102,204}
\definecolor{lightblue}{RGB}{220,235,252}
\definecolor{darkgreen}{RGB}{0,100,0}
\definecolor{lightgreen}{RGB}{220,255,220}
\definecolor{darkred}{RGB}{150,0,0}
\definecolor{lightred}{RGB}{255,230,230}
\definecolor{lightyellow}{RGB}{255,255,220}
\definecolor{orange}{RGB}{200,100,0}
\definecolor{purple}{RGB}{100,0,150}
\definecolor{lightgray}{RGB}{245,245,245}
\definecolor{darkgray}{RGB}{80,80,80}
\definecolor{teal}{RGB}{0,128,128}

% ─────────────────────────────────────────────
%  THEOREM ENVIRONMENTS
% ─────────────────────────────────────────────
\tcbuselibrary{most}

\newtcbtheorem[number within=section]{theorem}{Theorem}{
  enhanced, breakable,
  colback=lightblue, colframe=darkblue, fonttitle=\bfseries,
  attach boxed title to top left={yshift=-2mm,xshift=4mm},
  boxed title style={colback=darkblue,colframe=darkblue}
}{thm}

\newtcbtheorem[number within=section]{definition}{Definition}{
  enhanced, breakable,
  colback=lightgreen, colframe=darkgreen, fonttitle=\bfseries,
  attach boxed title to top left={yshift=-2mm,xshift=4mm},
  boxed title style={colback=darkgreen,colframe=darkgreen}
}{def}

\newtcbtheorem[number within=section]{example}{Example}{
  enhanced, breakable,
  colback=lightyellow, colframe=orange, fonttitle=\bfseries,
  attach boxed title to top left={yshift=-2mm,xshift=4mm},
  boxed title style={colback=orange,colframe=orange}
}{ex}

\newtcbtheorem[number within=section]{exercise}{Exercise}{
  enhanced, breakable,
  colback=lightred, colframe=darkred, fonttitle=\bfseries,
  attach boxed title to top left={yshift=-2mm,xshift=4mm},
  boxed title style={colback=darkred,colframe=darkred}
}{exc}

\newtcolorbox{keyidea}{
  enhanced, breakable,
  colback=lightyellow, colframe=purple, fonttitle=\bfseries,
  title=Key Idea, attach boxed title to top left={yshift=-2mm,xshift=4mm},
  boxed title style={colback=purple,colframe=purple}
}

\newtcolorbox{intuition}{
  enhanced, breakable,
  colback=lightgray, colframe=teal, fonttitle=\bfseries,
  title=Intuition, attach boxed title to top left={yshift=-2mm,xshift=4mm},
  boxed title style={colback=teal,colframe=teal}
}

\newtcolorbox{warning}{
  enhanced, breakable,
  colback=lightred, colframe=darkred, fonttitle=\bfseries,
  title={$\blacktriangleright$ Common Pitfall}, 
  attach boxed title to top left={yshift=-2mm,xshift=4mm},
  boxed title style={colback=darkred,colframe=darkred}
}

% ─────────────────────────────────────────────
%  HEADERS / FOOTERS
% ─────────────────────────────────────────────
\pagestyle{fancy}
\fancyhf{}
\fancyhead[L]{\small\textcolor{darkblue}{\textbf{Cook--Levin Theorem}}}
\fancyhead[R]{\small\textcolor{darkblue}{Theory of Computation}}
\fancyfoot[C]{\small\textcolor{darkgray}{\thepage}}
\renewcommand{\headrulewidth}{0.6pt}
\renewcommand{\footrulewidth}{0.4pt}

% ─────────────────────────────────────────────
%  SECTION STYLING
% ─────────────────────────────────────────────
\titleformat{\section}[block]
  {\Large\bfseries\color{darkblue}}
  {\thesection.}{0.5em}{}
  [\vspace{-0.3em}\textcolor{midblue}{\rule{\linewidth}{1pt}}]

\titleformat{\subsection}[block]
  {\large\bfseries\color{darkblue!80}}
  {\thesubsection.}{0.5em}{}

\titleformat{\subsubsection}[block]
  {\normalsize\bfseries\color{darkblue!60}}
  {\thesubsubsection.}{0.5em}{}

% ─────────────────────────────────────────────
%  MATH MACROS
% ─────────────────────────────────────────────
\newcommand{\P}{\mathbf{P}}
\newcommand{\NP}{\mathbf{NP}}
\newcommand{\SAT}{\textsc{Sat}}
\newcommand{\TSAT}{\textsc{3-Sat}}
\newcommand{\CLIQUE}{\textsc{Clique}}
\newcommand{\VERTEX}{\textsc{Vertex-Cover}}
\newcommand{\HAMPATH}{\textsc{Ham-Path}}
\newcommand{\SUBSET}{\textsc{Subset-Sum}}
\newcommand{\enc}[1]{\langle #1 \rangle}
\newcommand{\poly}{\mathrm{poly}}
\newcommand{\leqp}{\leq_p}

% ─────────────────────────────────────────────
%  DOCUMENT
% ─────────────────────────────────────────────
\begin{document}

% ══════════════════════════════════════════════
%  TITLE PAGE
% ══════════════════════════════════════════════
\begin{titlepage}
\begin{tikzpicture}[remember picture, overlay]
  % Top banner
  \fill[darkblue] (current page.north west) rectangle +(\paperwidth, -3.5cm);
  % Bottom accent
  \fill[midblue] (current page.south west) rectangle +(\paperwidth, 1.2cm);
  % Side stripe
  \fill[midblue!60] ([xshift=13.5cm]current page.south west) 
                    rectangle ++(1cm, \paperheight);
\end{tikzpicture}

\vspace*{2cm}
\begin{center}
  {\color{white}\Huge\bfseries Cook--Levin Theorem}\\[0.6em]
  {\color{white!80}\Large A Comprehensive Undergraduate Lecture}
\end{center}

\vspace{2cm}
\begin{center}
\begin{tikzpicture}[scale=0.85]
  % NP bubble
  \fill[midblue!30,rounded corners=18pt] (-3.5,-3) rectangle (3.5,3);
  \node[font=\large\bfseries,color=darkblue] at (0,2.2) {\NP};
  % P bubble inside
  \fill[lightblue,rounded corners=10pt] (-1.5,-1.5) rectangle (1.5,1.5);
  \node[font=\bfseries,color=darkblue] at (0,0.6) {\P};
  % SAT star
  \node[draw=darkred,fill=lightred,star,star points=5,star point ratio=2.2,
        minimum size=1.2cm,font=\small\bfseries,color=darkred] at (0,-0.3) {\SAT};
  % Arrows from other problems to SAT
  \node[font=\small\bfseries,color=purple] at (-2.5, 0.5) {\CLIQUE};
  \draw[-{Stealth[scale=1.2]},purple,thick] (-1.8,0.5) -- (-0.8,-0.1);
  \node[font=\small\bfseries,color=darkgreen] at (2.5, 0.5) {\HAMPATH};
  \draw[-{Stealth[scale=1.2]},darkgreen,thick] (1.75,0.5) -- (0.8,-0.1);
  \node[font=\small\bfseries,color=orange] at (0,-2.2) {\SUBSET};
  \draw[-{Stealth[scale=1.2]},orange,thick] (0,-1.75) -- (0,-0.75);
  % label
  \node[font=\itshape\small,color=darkgray] at (0,-3.5) 
    {Every \NP\ problem reduces to \SAT};
\end{tikzpicture}
\end{center}

\vspace{1.5cm}
\begin{center}
  \textcolor{darkgray}{\large Theory of Computation $\cdot$ Complexity Theory}\\[1em]
  \textcolor{darkgray}{\normalsize 
  \begin{tabular}{ll}
    \textbf{Topics:} & Turing Machines, Complexity Classes, Polynomial Reduction,\\
    & Boolean Satisfiability, NP-Completeness\\[0.4em]
    \textbf{Level:} & Undergraduate (2nd--3rd Year)\\[0.4em]
    \textbf{Prerequisites:} & Automata Theory, Basic Logic, Algorithms\\
  \end{tabular}
  }
\end{center}
\end{titlepage}

% ══════════════════════════════════════════════
%  TABLE OF CONTENTS
% ══════════════════════════════════════════════
\tableofcontents
\newpage

% ══════════════════════════════════════════════
\section{Motivation and Big Picture}
% ══════════════════════════════════════════════

\begin{intuition}
Imagine you have thousands of hard computational problems---scheduling flights, 
packing bags optimally, finding the shortest route through 50 cities. 
Is each problem \emph{independently} hard, or is there some deep structure 
they all share? The Cook--Levin theorem reveals a stunning answer:
\textbf{there is one single problem to which all of them reduce.}
\end{intuition}

\subsection{The Central Question of Computer Science}

The most famous open problem in all of mathematics and computer science is:

\begin{center}
\begin{tikzpicture}
  \node[draw=darkblue,fill=lightblue,rounded corners=8pt,
        inner sep=14pt,font=\Large\bfseries,color=darkblue,
        drop shadow] 
    {\textbf{Does $\P = \NP$?}};
\end{tikzpicture}
\end{center}

\medskip
In plain English:
\begin{itemize}[leftmargin=2em]
  \item \textbf{P} = problems whose solutions can be \emph{found} quickly (in polynomial time).
  \item \textbf{NP} = problems whose solutions can be \emph{verified} quickly (in polynomial time).
  \item The question asks: does ``verifiable quickly'' imply ``solvable quickly''?
\end{itemize}

\subsection{Why the Cook--Levin Theorem Matters}

The Cook--Levin theorem (1971--1973) does two things:
\begin{enumerate}[leftmargin=2em]
  \item \textbf{Defines NP-completeness}: it gives us a formal notion of the ``hardest'' problems in NP.
  \item \textbf{Shows \SAT\ is the first NP-complete problem}: every NP problem can be translated into SAT in polynomial time.
\end{enumerate}

If someone found a polynomial-time algorithm for SAT, \emph{every} NP problem 
would become tractable overnight. This is why SAT is called the ``master'' 
or ``canonical'' hard problem.

\bigskip
\begin{center}
\begin{tikzpicture}[node distance=1.5cm, every node/.style={font=\small}]
  \node[draw=darkblue,fill=lightblue,rounded corners,minimum width=3cm,
        minimum height=0.8cm,align=center] (SAT) {\textbf{SAT}\\(NP-complete)};
  \foreach \prob/\angle/\col in {
      Clique/135/purple,
      {3-Coloring}/100/teal,
      {Vertex Cover}/65/darkgreen,
      {Hamiltonian Path}/30/orange,
      {Subset Sum}/330/darkred,
      {Traveling Salesman}/295/midblue,
      {Integer Programming}/250/brown,
      {Circuit SAT}/215/darkgray}{
    \node[draw=\col,fill=\col!15,rounded corners,
          minimum width=2.3cm,minimum height=0.7cm,align=center,
          xshift={3.5cm*cos(\angle)},yshift={3.5cm*sin(\angle)}]
      (\prob) {\prob};
    \draw[-{Stealth},\col,thick] (\prob) -- (SAT);
  }
  \node[below=0.4cm of SAT,font=\itshape\small,color=darkgray] 
    {All arrows are polynomial-time reductions};
\end{tikzpicture}
\end{center}

\newpage

% ══════════════════════════════════════════════
\section{Prerequisite Background}
% ══════════════════════════════════════════════

\subsection{Turing Machines: The Universal Computational Model}

\begin{definition}{Turing Machine (TM)}{tm}
A \textbf{Turing Machine} $M$ is a 7-tuple 
$M = (Q, \Sigma, \Gamma, \delta, q_0, q_\mathrm{accept}, q_\mathrm{reject})$ where:
\begin{itemize}[leftmargin=2em]
  \item $Q$ = finite set of states
  \item $\Sigma$ = input alphabet ($\sqcup \notin \Sigma$)
  \item $\Gamma \supseteq \Sigma \cup \{\sqcup\}$ = tape alphabet ($\sqcup$ is blank)
  \item $\delta: Q \times \Gamma \to Q \times \Gamma \times \{L, R\}$ = transition function
  \item $q_0 \in Q$ = start state
  \item $q_\mathrm{accept}, q_\mathrm{reject} \in Q$ = halting states
\end{itemize}
\end{definition}

\begin{figure}[H]
\centering
\begin{tikzpicture}[scale=1]
  % Tape cells
  \foreach \i/\sym in {0/1,1/0,2/1,3/1,4/\sqcup,5/\sqcup,6/\sqcup}{
    \draw[thick,darkblue] (\i*1.1,0) rectangle ++(1.1,0.8);
    \node[font=\large] at (\i*1.1+0.55,0.4) {$\sym$};
  }
  \node[font=\small,color=darkgray] at (-0.5,0.4) {\ldots};
  \node[font=\small,color=darkgray] at (7.8,0.4) {\ldots};
  
  % Read/Write head
  \draw[darkred,very thick,-{Stealth[scale=1.5]}] (2.75,-0.5) -- (2.75,0);
  \node[draw=darkred,fill=lightred,rounded corners=4pt,
        minimum width=1.5cm,minimum height=0.6cm,
        font=\small\bfseries,color=darkred] at (2.75,-1.1) {Head ($q_i$)};
  
  % Labels
  \node[font=\small\itshape,color=darkgray] at (3.85,1.2) {Infinite Tape};
  \draw[darkgray,->] (3.2,1.1) -- (1.5,0.9);
  \draw[darkgray,->] (4.2,1.1) -- (5.5,0.9);
  
  % Control unit
  \node[draw=darkblue,fill=lightblue,rounded corners=8pt,
        minimum width=2.8cm,minimum height=1cm,
        font=\small\bfseries,align=center] at (8.5,-0.2)
    {Finite\\Control Unit};
  \draw[darkblue,<->,thick] (7.1,-0.2) -- (6.1,-0.2);
  \node[font=\tiny,color=darkgray] at (6.6,0.1) {state};
\end{tikzpicture}
\caption{Structure of a Turing Machine: an infinite tape, a read/write head, and a finite control.}
\end{figure}

\subsection{Computational Complexity Classes}

\begin{definition}{The Class P}{class_p}
$\P$ is the class of decision problems (yes/no problems) solvable by a 
\emph{deterministic} Turing machine in time $O(n^k)$ for some constant $k$, 
where $n$ is the input size.
\[
  \P = \bigcup_{k \geq 1} \mathbf{DTIME}(n^k)
\]
\end{definition}

\begin{example}{Problems in P}{p_examples}
\begin{itemize}[leftmargin=2em]
  \item Sorting $n$ numbers: $O(n \log n)$ --- in \P
  \item Shortest path (Dijkstra): $O(n^2)$ --- in \P
  \item Primality testing (AKS algorithm): $O(\log^{6} n)$ --- in \P
  \item Matrix multiplication: $O(n^{2.37})$ --- in \P
\end{itemize}
\end{example}

\begin{definition}{The Class NP}{class_np}
$\NP$ is the class of decision problems where every ``yes'' instance has a 
\emph{polynomial-length certificate} (proof/witness) that can be verified 
in polynomial time by a deterministic TM.
\[
  L \in \NP \iff \exists \text{ verifier } V \text{ s.t. } 
  \forall x \in L,\; \exists w,\; |w| \leq |x|^k,\; V(x,w) \text{ accepts in poly time}
\]
\end{definition}

\begin{intuition}
Think of \NP\ problems like a \textbf{crossword puzzle}:
\begin{itemize}[leftmargin=2em]
  \item \textbf{Solving} it from scratch might take a very long time.
  \item \textbf{Checking} someone else's completed solution is easy---you just verify each word fits.
\end{itemize}
\end{intuition}

\begin{example}{Problems in NP}{np_examples}
\begin{center}
\begin{tabular}{>{\bfseries}p{3.5cm} p{4.5cm} p{4cm}}
\toprule
Problem & Question & Certificate \\
\midrule
Graph Coloring & Can graph $G$ be 3-colored? & The coloring assignment \\
Hamiltonian Path & Does $G$ have a Ham.\ path? & The path itself \\
Subset Sum & Does a subset sum to $t$? & The subset \\
SAT & Is formula $\phi$ satisfiable? & The satisfying assignment \\
Clique & Does $G$ have a $k$-clique? & The $k$ vertices \\
\bottomrule
\end{tabular}
\end{center}
\end{example}

\begin{figure}[H]
\centering
\begin{tikzpicture}[scale=1.1]
  % EXPTIME
  \fill[gray!20,rounded corners=20pt] (-5,-3.3) rectangle (5,3.3);
  \node[font=\bfseries\small,color=darkgray] at (0,3.0) {\textbf{EXPTIME}};
  % NP
  \fill[midblue!25,rounded corners=14pt] (-3.8,-2.2) rectangle (3.8,2.2);
  \node[font=\bfseries,color=darkblue] at (2.8,1.8) {\NP};
  % co-NP  
  \fill[orange!25,rounded corners=14pt] (-3.6,-2.0) rectangle (3.6,-0.2);
  \node[font=\bfseries\small,color=orange] at (2.5,-0.6) {co-\NP};
  % P
  \fill[lightblue,rounded corners=10pt] (-2,-1.5) rectangle (2,1.5);
  \node[font=\bfseries,color=darkblue] at (0,0) {\Large$\P$};
  % labels
  \node[font=\small,color=purple,align=center] at (-3.0,1.5) 
    {NP-complete\\problems};
  \draw[purple,->] (-2.2,1.4) -- (-0.5,0.7);
\end{tikzpicture}
\caption{The complexity class landscape. Whether $\P = \NP$, whether $\NP = \text{co-NP}$, etc.,\ are all open questions.}
\end{figure}

\subsection{Boolean Logic: Building Blocks of SAT}

\begin{definition}{Boolean Formula}{bool_formula}
A \textbf{Boolean formula} is built from:
\begin{itemize}[leftmargin=2em]
  \item \textbf{Variables}: $x_1, x_2, \ldots, x_n \in \{0,1\}$ (false/true)
  \item \textbf{Literals}: $x_i$ or $\neg x_i$
  \item \textbf{Connectives}: AND ($\land$), OR ($\lor$), NOT ($\neg$)
\end{itemize}
A \textbf{clause} is a disjunction (OR) of literals: $(x_1 \lor \neg x_2 \lor x_3)$.\\
A formula is in \textbf{Conjunctive Normal Form (CNF)} if it is an AND of clauses.
\end{definition}

\begin{example}{CNF Formula}{cnf_ex}
\[
  \phi = \underbrace{(x_1 \lor \neg x_2)}_{\text{clause 1}} 
       \;\land\; 
       \underbrace{(\neg x_1 \lor x_3 \lor x_2)}_{\text{clause 2}} 
       \;\land\;
       \underbrace{(\neg x_3)}_{\text{clause 3}}
\]
Assignment $x_1 = 1,\, x_2 = 1,\, x_3 = 0$: evaluate each clause:
\begin{itemize}[leftmargin=2em]
  \item Clause 1: $(1 \lor 0) = 1$ \checkmark
  \item Clause 2: $(0 \lor 0 \lor 1) = 1$ \checkmark
  \item Clause 3: $(1) = 1$ \checkmark
\end{itemize}
So $\phi$ is \textbf{satisfiable} (SAT).
\end{example}

\newpage

% ══════════════════════════════════════════════
\section{Boolean Satisfiability (SAT)}
% ══════════════════════════════════════════════

\begin{definition}{SAT Problem}{sat_def}
\[
  \SAT = \{ \enc{\phi} \mid \phi \text{ is a satisfiable Boolean formula in CNF} \}
\]
\textbf{Input:} A Boolean formula $\phi$ in CNF over variables $x_1, \ldots, x_n$.\\
\textbf{Question:} Does there exist an assignment $\alpha: \{x_1,\ldots,x_n\} \to \{0,1\}$ 
such that $\phi(\alpha) = 1$?
\end{definition}

\subsection{Naive Algorithm for SAT}

The brute-force approach tries all $2^n$ assignments:

\begin{algorithm}[H]
\caption{Brute-Force SAT Solver}
\begin{algorithmic}[1]
\Require CNF formula $\phi$ with variables $x_1, \ldots, x_n$
\Ensure \textbf{SAT} if satisfiable, \textbf{UNSAT} otherwise
\For{each assignment $\alpha \in \{0,1\}^n$}   \Comment{$2^n$ iterations}
  \If{$\phi(\alpha) = 1$}
    \State \Return \textbf{SAT} (with witness $\alpha$)
  \EndIf
\EndFor
\State \Return \textbf{UNSAT}
\end{algorithmic}
\end{algorithm}

\medskip
This runs in $O(2^n \cdot |\phi|)$ time --- \textbf{exponential}. No polynomial algorithm is known!

\subsection{SAT is in NP}

\begin{theorem}{SAT $\in$ NP}{sat_in_np}
The Boolean Satisfiability problem is in NP.
\end{theorem}

\begin{proof}
Given a formula $\phi$ and a proposed certificate (assignment) $\alpha$:
\begin{enumerate}[leftmargin=2em]
  \item The certificate $\alpha$ has length $n$ bits --- polynomial in $|\phi|$.
  \item Evaluating $\phi(\alpha)$: substitute each variable and check each clause.
    This takes $O(|\phi|)$ time --- polynomial.
\end{enumerate}
Therefore SAT has an efficient verifier, so SAT $\in$ NP. \qed
\end{proof}

\begin{figure}[H]
\centering
\begin{tikzpicture}[scale=0.95, node distance=1.4cm and 2cm,
  box/.style={draw,rounded corners=5pt,minimum height=0.9cm,
              minimum width=2.5cm,align=center,font=\small}]
  % Input
  \node[box,fill=lightblue,draw=darkblue] (input) 
    {Input $\phi$\\(CNF formula)};
  % Certificate  
  \node[box,fill=lightgreen,draw=darkgreen,below=of input] (cert) 
    {Certificate $\alpha$\\(assignment, $n$ bits)};
  % Verifier
  \node[box,fill=lightyellow,draw=orange,right=2.5cm of input,
        yshift=-0.7cm,minimum height=2.2cm] (verifier) 
    {Poly-Time\\Verifier $V$};
  % Output yes/no
  \node[box,fill=lightred,draw=darkred,right=2cm of verifier] (out)
    {Accept / Reject};
  
  \draw[-{Stealth},thick,darkblue] (input) -- (verifier);
  \draw[-{Stealth},thick,darkgreen] (cert) -- (verifier);
  \draw[-{Stealth},thick,darkred] (verifier) -- (out);
  
  % Annotation
  \node[font=\small\itshape,color=darkgray,below=0.2cm of verifier]
    {runs in $O(|\phi|)$ time};
\end{tikzpicture}
\caption{The NP verification procedure for SAT.}
\end{figure}

\newpage

% ══════════════════════════════════════════════
\section{Polynomial-Time Reductions}
% ══════════════════════════════════════════════

\subsection{What is a Reduction?}

\begin{definition}{Polynomial-Time Reduction}{poly_reduce}
A \textbf{polynomial-time many-one reduction} from language $A$ to language $B$, 
written $A \leqp B$, is a function $f: \Sigma^* \to \Sigma^*$ such that:
\begin{enumerate}[leftmargin=2em]
  \item $f$ is computable in polynomial time.
  \item For all $x$: $x \in A \iff f(x) \in B$.
\end{enumerate}
\end{definition}

\begin{intuition}
A reduction $A \leqp B$ says: ``If I could solve $B$ quickly, I could solve $A$ quickly.''
Equivalently: $B$ is \emph{at least as hard} as $A$.

\medskip
Think of it like a \textbf{translation service}: any instance of problem $A$ can be 
``translated'' into an instance of problem $B$ in polynomial time, without changing 
the yes/no answer.
\end{intuition}

\begin{figure}[H]
\centering
\begin{tikzpicture}[scale=0.95]
  % Problem A box
  \node[draw=darkblue,fill=lightblue,rounded corners,
        minimum width=2.5cm,minimum height=1.5cm,align=center,
        font=\bfseries] (A) at (0,0) {Problem $A$\\instance $x$};
  % Arrow with reduction function
  \draw[-{Stealth[scale=1.5]},midblue,very thick] (2,0) -- (5,0)
    node[midway,above,font=\small\bfseries,color=midblue] {$f$ (poly time)}
    node[midway,below,font=\small\itshape,color=darkgray] {transformation};
  % Problem B box
  \node[draw=darkred,fill=lightred,rounded corners,
        minimum width=2.5cm,minimum height=1.5cm,align=center,
        font=\bfseries] (B) at (6.5,0) {Problem $B$\\instance $f(x)$};
  % Oracle for B
  \node[draw=purple,fill=purple!15,rounded corners,
        minimum width=2cm,minimum height=1cm,
        font=\small\bfseries] (oracle) at (10,0) {Oracle\\for $B$};
  \draw[-{Stealth},thick,purple] (B) -- (oracle);
  % Answer
  \node[font=\bfseries] (ans) at (12.5,0) {Yes/No};
  \draw[-{Stealth},thick] (oracle) -- (ans);
  % Answer flows back to A
  \draw[-{Stealth},thick,darkgray,dashed,bend left=30] 
    (12.5,-0.5) to node[midway,below,font=\small\itshape] 
    {same answer for $A$} (0,-0.5);
  
  % Condition
  \node[font=\small\bfseries,color=darkgreen] at (6.5,-2) 
    {$x \in A \iff f(x) \in B$};
\end{tikzpicture}
\caption{Polynomial-time reduction: an oracle for $B$ gives us a solver for $A$.}
\end{figure}

\subsection{Properties of Reductions}

\begin{theorem}{Transitivity of Reductions}{trans}
If $A \leqp B$ and $B \leqp C$, then $A \leqp C$.
\end{theorem}

\begin{proof}
Let $f$ be the poly-time reduction from $A$ to $B$, and $g$ the poly-time reduction from $B$ to $C$.
Define $h(x) = g(f(x))$.
\begin{itemize}[leftmargin=2em]
  \item $h$ is computable in polynomial time (composition of two poly-time functions).
  \item $x \in A \iff f(x) \in B \iff g(f(x)) \in C$. \qed
\end{itemize}
\end{proof}

\begin{theorem}{Reduction Transfers Tractability}{tractability}
If $A \leqp B$ and $B \in \P$, then $A \in \P$.
\end{theorem}

\begin{proof}
Given input $x$ for $A$:
\begin{enumerate}[leftmargin=2em]
  \item Compute $f(x)$ in poly time $p_f(|x|)$.
  \item Run the poly-time algorithm for $B$ on $f(x)$. Time: $p_B(|f(x)|) \leq p_B(p_f(|x|))$.
\end{enumerate}
The composition is a polynomial (poly of poly), so $A \in \P$. \qed
\end{proof}

\subsection{NP-Hardness and NP-Completeness}

\begin{definition}{NP-Hard and NP-Complete}{npc}
\begin{itemize}[leftmargin=2em]
  \item A problem $B$ is \textbf{NP-hard} if for every $A \in \NP$: $A \leqp B$.
  \item A problem $B$ is \textbf{NP-complete} if $B \in \NP$ \emph{and} $B$ is NP-hard.
\end{itemize}
\end{definition}

\begin{keyidea}
\textbf{NP-complete problems are the hardest problems in NP.} If any NP-complete problem 
is in P, then $\P = \NP$. If any NP-complete problem is not in P, then $\P \neq \NP$.
\end{keyidea}

\newpage

% ══════════════════════════════════════════════
\section{The Cook--Levin Theorem}
% ══════════════════════════════════════════════

\subsection{Statement}

\begin{theorem}{Cook--Levin Theorem (1971, 1973)}{cook_levin}
\textbf{SAT is NP-complete.} That is:
\begin{enumerate}[leftmargin=2em]
  \item $\SAT \in \NP$, and
  \item For every $L \in \NP$: $L \leqp \SAT$.
\end{enumerate}
\end{theorem}

\begin{center}
\begin{tikzpicture}
  \node[draw=darkblue,fill=lightblue,rounded corners=10pt,
        inner sep=16pt,align=center,font=\large,
        drop shadow={shadow xshift=2pt,shadow yshift=-2pt}]
    {\textbf{SAT is the first NP-complete problem.}\\[4pt]
     Every problem in NP ``hides inside'' SAT.};
\end{tikzpicture}
\end{center}

\subsection{Historical Context}

\begin{center}
\begin{tikzpicture}[scale=0.9]
  % Timeline
  \draw[darkblue,very thick,->] (0,0) -- (12,0);
  % Ticks
  \foreach \x/\yr/\lbl/\col in {
    1/1936/{Turing Machine\\invented (Turing)}/darkgray,
    3.5/1965/{\P\ class\\defined}/darkblue,
    6/1971/{Cook proves\\SAT NP-complete}/darkred,
    8.5/1972/{Karp's 21 NP-\\complete problems}/purple,
    10.5/1973/{Levin's\\independent proof}/teal}{
    \draw[\col,thick] (\x,-0.2) -- (\x,0.2);
    \node[\col,above,font=\tiny\bfseries,align=center,yshift=2pt] at (\x,0.2) {\yr};
    \node[\col,below,font=\tiny,align=center,yshift=-2pt] at (\x,-0.2) {\lbl};
  }
\end{tikzpicture}
\end{center}

\medskip
\begin{itemize}[leftmargin=2em]
  \item \textbf{Stephen Cook} (University of Toronto, 1971): proved SAT is NP-complete in his landmark paper.
  \item \textbf{Richard Karp} (UC Berkeley, 1972): reduced 21 major combinatorial problems to SAT, founding NP-completeness theory.
  \item \textbf{Leonid Levin} (Soviet Union, 1973): independently proved essentially the same result.
\end{itemize}

\subsection{High-Level Proof Idea}

\begin{intuition}
Any NP problem $L$ has a polynomial-time verifier $V$. Given input $x$:
\begin{enumerate}[leftmargin=2em]
  \item The verifier $V(x, w)$ runs on a Turing machine.
  \item A TM computation is a sequence of configurations (``snapshots'').
  \item We can encode ``does $V$ accept?'' as a Boolean formula $\phi_{x}$.
  \item $V$ accepts iff $\phi_x$ is satisfiable --- with the satisfying assignment encoding the certificate $w$.
\end{enumerate}
\end{intuition}

\begin{figure}[H]
\centering
\begin{tikzpicture}[scale=0.9, node distance=1.5cm and 0.5cm]
  \node[draw=darkblue,fill=lightblue,rounded corners,minimum width=3cm,
        minimum height=0.9cm,align=center,font=\bfseries] (prob) 
    {$L \in \NP$\\instance $x$};
  \node[draw=darkgreen,fill=lightgreen,rounded corners,minimum width=3cm,
        minimum height=0.9cm,align=center,font=\bfseries,right=2cm of prob] (ver)
    {Poly-Time\\Verifier $V$};
  \node[draw=purple,fill=purple!15,rounded corners,minimum width=3cm,
        minimum height=0.9cm,align=center,font=\bfseries,right=2cm of ver] (tm)
    {TM Tableau\\(computation table)};
  \node[draw=darkred,fill=lightred,rounded corners,minimum width=3cm,
        minimum height=0.9cm,align=center,font=\bfseries,right=2cm of tm] (phi)
    {CNF Formula $\phi_x$};
  
  \draw[-{Stealth},thick,darkblue] (prob) -- (ver) 
    node[midway,above,font=\tiny] {has};
  \draw[-{Stealth},thick,darkgreen] (ver) -- (tm) 
    node[midway,above,font=\tiny] {encoded as};
  \draw[-{Stealth},thick,purple] (tm) -- (phi)
    node[midway,above,font=\tiny] {encoded as};
  
  \node[font=\small\itshape,color=darkgray,below=0.4cm of ver] 
    {``$V(x,w)$ accepts''};
  \node[font=\small\itshape,color=darkgray,below=0.4cm of phi]
    {satisfiable $\iff x \in L$};
\end{tikzpicture}
\caption{The chain of encodings in the Cook--Levin proof.}
\end{figure}

\newpage
\subsection{Detailed Proof of Cook--Levin}

We prove that for any $L \in \NP$, there is a poly-time reduction $L \leqp \SAT$.

\subsubsection{Step 1: The Verifier}

Since $L \in \NP$, there exists a poly-time verifier $V$ and a polynomial $p$ such that:
\[
  x \in L \iff \exists w,\; |w| \leq p(|x|),\; V(x,w) \text{ accepts}
\]
Say $V$ runs in time $\leq q(n)$ for polynomial $q$, where $n = |x|$.

\subsubsection{Step 2: The Tableau}

A TM \textbf{tableau} is a grid recording the entire computation:

\begin{center}
\begin{tikzpicture}[scale=0.8]
  \def\rows{5}
  \def\cols{8}
  % Grid
  \foreach \r in {0,...,\rows}{
    \foreach \c in {0,...,\cols}{
      \draw[darkblue!30] (\c*1.1, \r*0.8) rectangle ++(1.1,0.8);
    }
  }
  % Labels
  \node[rotate=90,font=\small\bfseries,color=darkblue] at (-0.5, 2.4) {Time steps};
  \node[font=\small\bfseries,color=darkblue] at (4.95, -0.5) {Tape cells};
  
  % Row labels
  \foreach \r/\lbl in {0/t=0,1/t=1,2/t=2,3/\vdots,4/{t=q(n)}}{
    \node[font=\small,color=darkgray] at (-1.0, \r*0.8+0.4) {\lbl};
  }
  % Initial config
  \foreach \c/\sym in {0/q_0,1/x_1,2/x_2,3/\cdots,4/x_n,5/w_1,6/\cdots,7/\sqcup}{
    \node[font=\scriptsize] at (\c*1.1+0.55, 0.4) {$\sym$};
  }
  % Accept state
  \node[draw=darkgreen,fill=lightgreen,rounded corners=3pt,
        font=\scriptsize\bfseries,minimum width=1.0cm] 
    at (2.75, 4.4) {$q_\mathrm{acc}$};
  
  % Highlight head movement
  \draw[darkred,very thick,rounded corners=4pt] 
    (0,0) rectangle (1.1,0.8);
  \draw[darkred,very thick,rounded corners=4pt] 
    (1.1,0.8) rectangle (2.2,1.6);
  \draw[darkred,very thick,rounded corners=4pt] 
    (2.2,1.6) rectangle (3.3,2.4);
  
  % Legend
  \node[font=\tiny\itshape,color=darkred] at (10,1.2) 
    {red box = head position};
\end{tikzpicture}
\end{center}

The tableau has dimensions $q(n) \times q(n)$ (time steps $\times$ tape width).
Each cell $(i,j)$ contains a symbol from $\Gamma \cup Q$ (the tape symbol and possibly the state if the head is at column $j$ at time $i$).

\subsubsection{Step 3: Variables}

For each cell $(i,j)$ with $0 \leq i \leq q(n)$ and $0 \leq j \leq q(n)$, 
and for each symbol $s \in C = \Gamma \cup (Q \times \Gamma)$ (an alphabet of size $|C|$):
\[
  x_{i,j,s} = 
  \begin{cases} 
    1 & \text{if cell $(i,j)$ of the tableau contains symbol $s$} \\
    0 & \text{otherwise}
  \end{cases}
\]
Total variables: $O(q(n)^2 \cdot |C|) = O(n^{2k})$ --- polynomial.

\subsubsection{Step 4: The Four Groups of Clauses}

We write the formula $\phi_x = \phi_{\text{cell}} \land \phi_{\text{start}} \land \phi_{\text{move}} \land \phi_{\text{accept}}$.

\medskip
\noindent\textbf{(a) Cell Constraints} $\phi_{\text{cell}}$: Each cell contains exactly one symbol.
\[
  \phi_{\text{cell}} = \bigwedge_{i,j} \left( 
    \underbrace{\bigvee_{s \in C} x_{i,j,s}}_{\text{at least one}} 
    \;\land\; 
    \underbrace{\bigwedge_{s \neq t} (\neg x_{i,j,s} \lor \neg x_{i,j,t})}_{\text{at most one}}
  \right)
\]

\medskip
\noindent\textbf{(b) Start Constraints} $\phi_{\text{start}}$: The first row matches initial config.
\[
  \phi_{\text{start}} = x_{0,0,q_0} \land x_{0,1,x_1} \land \cdots \land x_{0,n,x_n} \land 
  x_{0,n+1,w_1} \land \cdots \land x_{0,n+p(n),w_{p(n)}} \land \cdots
\]
(The $w_i$ variables become the ``free'' variables encoding the certificate.)

\medskip
\noindent\textbf{(c) Move Constraints} $\phi_{\text{move}}$: Each row follows from the previous via $\delta$.
For every position $(i,j)$, a $2 \times 3$ window of cells must be legal:

\begin{center}
\begin{tikzpicture}[scale=0.85]
  % Window
  \foreach \c in {0,1,2}{
    \draw[darkblue,thick] (\c*1.2,0.8) rectangle ++(1.2,0.8);
    \draw[darkred,thick] (\c*1.2,0) rectangle ++(1.2,0.8);
  }
  \node[font=\small] at (0.6,1.2) {$j{-}1$};
  \node[font=\small] at (1.8,1.2) {$j$};
  \node[font=\small] at (3.0,1.2) {$j{+}1$};
  \node[rotate=90,font=\small,color=darkblue] at (-0.5,1.2) {row $i$};
  \node[rotate=90,font=\small,color=darkred] at (-0.5,0.4) {row $i+1$};
  
  % Arrow
  \draw[-{Stealth},thick] (4.5,0.9) -- (6,0.9) 
    node[midway,above,font=\tiny] {legal by $\delta$};
  
  % Legal window notation
  \node[font=\small] at (7.5,0.9) 
    {$\in \mathcal{W}$ (set of legal windows)};
\end{tikzpicture}
\end{center}

$\phi_{\text{move}} = \bigwedge_{i,j} \bigvee_{(a,b,c,a',b',c') \in \mathcal{W}} 
  (x_{i,j-1,a} \land x_{i,j,b} \land x_{i,j+1,c} \land x_{i+1,j-1,a'} \land x_{i+1,j,b'} \land x_{i+1,j+1,c'})$

Each 6-literal conjunction can be converted to CNF in $O(1)$ clauses (standard conversion).

\medskip
\noindent\textbf{(d) Accept Constraints} $\phi_{\text{accept}}$: The accepting state appears.
\[
  \phi_{\text{accept}} = \bigvee_{0 \leq j \leq q(n)} x_{q(n), j, q_\mathrm{accept}}
\]

\subsubsection{Step 5: Correctness}

\begin{theorem}{Correctness of the Reduction}{correct}
$x \in L \iff \phi_x$ is satisfiable.
\end{theorem}

\begin{proof}
\textbf{($\Rightarrow$)} If $x \in L$, then $\exists w$ such that $V(x,w)$ accepts in $\leq q(n)$ steps.
The resulting tableau is a valid accepting computation. Set $x_{i,j,s} = 1$ iff cell $(i,j)$ 
contains symbol $s$ in this tableau. This assignment satisfies all four groups of clauses.

\medskip
\noindent\textbf{($\Leftarrow$)} If $\phi_x$ is satisfiable with assignment $\alpha$:
\begin{itemize}[leftmargin=2em]
  \item $\phi_{\text{cell}}$ ensures each cell has exactly one symbol.
  \item $\phi_{\text{start}}$ ensures the first row is the correct start config.
  \item $\phi_{\text{move}}$ ensures each row is a valid TM transition.
  \item $\phi_{\text{accept}}$ ensures the last row contains $q_\mathrm{accept}$.
\end{itemize}
So the tableau describes a valid accepting computation of $V(x, w)$ where 
$w$ is extracted from columns $n+1$ to $n+p(n)$ of row 0. Thus $x \in L$. \qed
\end{proof}

\subsubsection{Step 6: Polynomial Size}

\begin{itemize}[leftmargin=2em]
  \item \textbf{Variables:} $O(q(n)^2 \cdot |C|) = O(n^{2k})$
  \item \textbf{Clauses:} 
    \begin{itemize}
      \item $\phi_{\text{cell}}$: $O(q(n)^2 \cdot |C|^2) = O(n^{2k})$ clauses
      \item $\phi_{\text{start}}$: $O(q(n))$ unit clauses
      \item $\phi_{\text{move}}$: $O(q(n)^2 \cdot |\mathcal{W}| \cdot |C|)$ clauses --- polynomial since $|\mathcal{W}|$ and $|C|$ are constants
      \item $\phi_{\text{accept}}$: $O(q(n))$ clauses
    \end{itemize}
  \item \textbf{Construction time:} $O(n^{2k})$ --- polynomial.
\end{itemize}

\medskip
\begin{center}
\begin{tikzpicture}
  \node[draw=darkgreen,fill=lightgreen,rounded corners=8pt,
        inner sep=12pt,align=center,font=\normalsize]
    {$L \leqp \SAT$ via a polynomial-time computable function $x \mapsto \phi_x$\\
     with $|\phi_x| = O(n^{2k})$. \textbf{QED.}};
\end{tikzpicture}
\end{center}

\newpage

% ══════════════════════════════════════════════
\section{Worked Examples}
% ══════════════════════════════════════════════

\subsection{Example 1: Encoding a Tiny TM Computation as SAT}

\begin{example}{Encoding a 2-Step Computation}{tiny_enc}
Consider a TM $M$ that accepts $\{0\}$ (strings equal to ``0'').
States: $Q = \{q_0, q_\text{acc}, q_\text{rej}\}$.
Tape alphabet: $\Gamma = \{0, 1, \sqcup\}$.

\medskip
Transition (simplified): $\delta(q_0, 0) = (q_\text{acc}, 0, R)$.

\medskip
\noindent\textbf{Input:} $x = 0$ (so $n = 1$). Take time bound $q(n) = 2$.

\medskip
\noindent\textbf{Tableau} ($3$ rows $\times$ $3$ columns):
\begin{center}
\begin{tabular}{c|ccc}
\toprule
Time & Cell 0 & Cell 1 & Cell 2 \\
\midrule
$t=0$ & $(q_0, 0)$ & $\sqcup$ & $\sqcup$ \\
$t=1$ & $0$ & $(q_\text{acc}, \sqcup)$ & $\sqcup$ \\
$t=2$ & $0$ & $(q_\text{acc}, \sqcup)$ & $\sqcup$ \\
\bottomrule
\end{tabular}
\end{center}

\medskip
\noindent\textbf{Some variables (partial list):}
\begin{itemize}[leftmargin=2em]
  \item $x_{0,0,(q_0,0)} = 1$: at time 0, cell 0 has symbol $(q_0, 0)$
  \item $x_{0,1,\sqcup} = 1$: at time 0, cell 1 is blank
  \item $x_{1,1,(q_\text{acc},\sqcup)} = 1$: at time 1, head at cell 1 in accept state
\end{itemize}

\medskip
\noindent\textbf{The SAT formula} $\phi_x$ is satisfied by exactly the assignment 
corresponding to this tableau, confirming $x = 0$ is accepted.
\end{example}

\subsection{Example 2: A Simple SAT Instance from Graph Coloring}

\begin{example}{Reducing 3-Coloring to SAT}{coloring_sat}
\textbf{Problem:} Can we 3-color the triangle graph $K_3$?

\medskip
\noindent\textbf{Variables:} For each vertex $v \in \{A,B,C\}$ and color $c \in \{R,G,B\}$:
$x_{v,c}$ = ``vertex $v$ has color $c$''.

\medskip
\noindent\textbf{Clauses:}

\begin{enumerate}[leftmargin=2em]
  \item \textit{Each vertex has at least one color:}
    \[(x_{A,R} \lor x_{A,G} \lor x_{A,B}) \land (x_{B,R} \lor x_{B,G} \lor x_{B,B}) \land (x_{C,R} \lor x_{C,G} \lor x_{C,B})\]
  \item \textit{Each vertex has at most one color (no two colors simultaneously):}
    \[(\neg x_{A,R} \lor \neg x_{A,G}) \land (\neg x_{A,R} \lor \neg x_{A,B}) \land (\neg x_{A,G} \lor \neg x_{A,B}) \land \cdots\]
  \item \textit{Adjacent vertices have different colors (edge $AB$):}
    \[(\neg x_{A,R} \lor \neg x_{B,R}) \land (\neg x_{A,G} \lor \neg x_{B,G}) \land (\neg x_{A,B} \lor \neg x_{B,B})\]
    (Similarly for edges $BC$ and $AC$.)
\end{enumerate}

\medskip
\noindent\textbf{Satisfying assignment:} $x_{A,R}=1,\; x_{B,G}=1,\; x_{C,B}=1$ (all others 0).
This encodes the coloring $A \mapsto R,\; B \mapsto G,\; C \mapsto B$, which is valid since $K_3$ can be 3-colored.

\begin{center}
\begin{tikzpicture}[scale=0.8]
  \node[circle,draw=darkred,fill=red!60,minimum size=0.9cm,font=\bfseries] (A) at (0,1.5) {A};
  \node[circle,draw=darkgreen,fill=green!60,minimum size=0.9cm,font=\bfseries] (B) at (-1.3,-0.8) {B};
  \node[circle,draw=blue!80,fill=blue!40,minimum size=0.9cm,font=\bfseries\color{white}] (C) at (1.3,-0.8) {C};
  \draw[very thick] (A)--(B)--(C)--(A);
  \node[font=\small\itshape,color=darkgray] at (0,-1.8) {Valid 3-coloring of $K_3$};
\end{tikzpicture}
\end{center}
\end{example}

\subsection{Example 3: Unsatisfiable Formula}

\begin{example}{UNSAT Instance}{unsat_ex}
Consider:
\[
  \phi = (x_1) \land (\neg x_1)
\]
\begin{itemize}[leftmargin=2em]
  \item If $x_1 = 1$: clause 1 satisfied, clause 2 = $\neg 1 = 0$. \textbf{FAIL.}
  \item If $x_1 = 0$: clause 1 = $0$. \textbf{FAIL.}
\end{itemize}
$\phi$ is \textbf{unsatisfiable}. The corresponding decision problem instance would be a ``No'' instance.

\medskip
In the Cook--Levin reduction, $\phi_x$ is unsatisfiable iff $x \notin L$ (no accepting computation exists).
\end{example}

\newpage

% ══════════════════════════════════════════════
\section{Consequences and Further NP-Complete Problems}
% ══════════════════════════════════════════════

\subsection{The Cascade of NP-Complete Problems}

Once Cook--Levin gave us SAT, Karp (1972) proved 21 more problems NP-complete 
by chaining polynomial reductions:

\begin{figure}[H]
\centering
\begin{tikzpicture}[
  node distance=1.3cm and 1.8cm,
  box/.style={draw,rounded corners=5pt,minimum width=2.5cm,
              minimum height=0.7cm,align=center,font=\small\bfseries},
  arr/.style={-{Stealth},thick}
]
  \node[box,fill=lightred,draw=darkred] (sat) {\SAT};
  \node[box,fill=lightblue,draw=darkblue,below left=of sat] (3sat) {\TSAT};
  \node[box,fill=lightblue,draw=darkblue,below right=of sat] (csat) {\textsc{Circuit-SAT}};
  \node[box,fill=lightgreen,draw=darkgreen,below=of 3sat] (clique) {\CLIQUE};
  \node[box,fill=lightgreen,draw=darkgreen,below left=of clique] (vc) {\VERTEX};
  \node[box,fill=lightgreen,draw=darkgreen,below right=of clique] (is) {\textsc{Ind.-Set}};
  \node[box,fill=purple!20,draw=purple,below=of 3sat,xshift=-3cm] (hp) {\HAMPATH};
  \node[box,fill=purple!20,draw=purple,below=of hp] (hc) {\textsc{Ham-Cycle}};
  \node[box,fill=purple!20,draw=purple,below=of hc] (tsp) {\textsc{TSP}};
  \node[box,fill=orange!20,draw=orange,below right=of 3sat] (ss) {\SUBSET};
  \node[box,fill=orange!20,draw=orange,below=of ss] (part) {\textsc{Partition}};
  \node[box,fill=orange!20,draw=orange,below=of part] (knap) {\textsc{Knapsack}};

  \draw[arr,darkred] (sat) -- (3sat);
  \draw[arr,darkred] (sat) -- (csat);
  \draw[arr,darkblue] (3sat) -- (clique);
  \draw[arr,darkblue] (clique) -- (vc);
  \draw[arr,darkblue] (clique) -- (is);
  \draw[arr,darkblue] (3sat) -- (hp);
  \draw[arr,purple] (hp) -- (hc);
  \draw[arr,purple] (hc) -- (tsp);
  \draw[arr,darkblue] (3sat) -- (ss);
  \draw[arr,orange] (ss) -- (part);
  \draw[arr,orange] (part) -- (knap);
  
  \node[font=\tiny\itshape,color=darkgray] at (0,-7) 
    {Each arrow is a polynomial-time reduction};
\end{tikzpicture}
\caption{A reduction tree rooted at SAT, showing how NP-completeness propagates.}
\end{figure}

\subsection{3-SAT: The Most Useful Variant}

\begin{definition}{3-SAT}{tsat_def}
\[
  \TSAT = \{ \enc{\phi} \mid \phi \text{ is a satisfiable CNF formula with exactly 3 literals per clause} \}
\]
\end{definition}

\begin{theorem}{3-SAT is NP-complete}{tsat_npc}
$\TSAT$ is NP-complete, i.e., $\SAT \leqp \TSAT$.
\end{theorem}

\begin{proof}[Proof sketch]
Given a clause $(l_1 \lor l_2 \lor l_3 \lor \cdots \lor l_k)$ with $k \neq 3$:

\medskip
\noindent\textbf{Case $k < 3$:} Pad with a fresh variable $y$:
\begin{itemize}[leftmargin=2em]
  \item $k=1$: $(l_1) \to (l_1 \lor y \lor \neg y)$ --- always satisfies the extra literals, forces $l_1$.
  \item $k=2$: $(l_1 \lor l_2) \to (l_1 \lor l_2 \lor y) \land (l_1 \lor l_2 \lor \neg y)$.
\end{itemize}

\noindent\textbf{Case $k > 3$:} Introduce helper variables $y_1, \ldots, y_{k-3}$:
\[
  (l_1 \lor l_2 \lor y_1) \land (\neg y_1 \lor l_3 \lor y_2) \land \cdots 
  \land (\neg y_{k-4} \lor l_{k-2} \lor y_{k-3}) \land (\neg y_{k-3} \lor l_{k-1} \lor l_k)
\]
One can verify this is satisfiable iff the original clause is satisfiable. \qed
\end{proof}

\begin{example}{Long Clause to 3-SAT}{longclause}
Convert $(x_1 \lor x_2 \lor x_3 \lor x_4 \lor x_5)$ (5 literals) to 3-SAT:
\begin{align*}
  &(x_1 \lor x_2 \lor y_1) \\
  &(\neg y_1 \lor x_3 \lor y_2) \\
  &(\neg y_2 \lor x_4 \lor x_5)
\end{align*}
with fresh variables $y_1, y_2$. This is a 3-CNF formula equivalent to the original clause.
\end{example}

\newpage

% ══════════════════════════════════════════════
\section{Implications for Computer Science}
% ══════════════════════════════════════════════

\subsection{What NP-Completeness Tells Us}

\begin{center}
\begin{tikzpicture}[scale=0.95]
  \node[draw=darkblue,fill=lightblue,rounded corners=8pt,
        inner sep=10pt,align=center,font=\small] (center) 
    at (0,0) {\textbf{Problem is NP-complete}};
  
  % Consequence nodes
  \node[draw=darkred,fill=lightred,rounded corners=5pt,
        minimum width=3.5cm,minimum height=1.2cm,align=center,font=\small] 
    at (-4.5, 2) {\textbf{No known poly-time algorithm}\\ likely needs exponential time};
  \node[draw=darkgreen,fill=lightgreen,rounded corners=5pt,
        minimum width=3.5cm,minimum height=1.2cm,align=center,font=\small] 
    at (4.5, 2) {\textbf{Solutions verifiable in poly time}\\ certificates exist};
  \node[draw=purple,fill=purple!15,rounded corners=5pt,
        minimum width=3.5cm,minimum height=1.2cm,align=center,font=\small] 
    at (-4.5, -2) {\textbf{All NP problems reduce to it}\\ hardest in NP};
  \node[draw=orange,fill=orange!15,rounded corners=5pt,
        minimum width=3.5cm,minimum height=1.2cm,align=center,font=\small] 
    at (4.5, -2) {\textbf{Solving it solves all NP}\\ $\P = \NP$ if in P};
  
  \draw[-{Stealth},thick,darkred] (-1.5,0.3) -- (-3,1.5);
  \draw[-{Stealth},thick,darkgreen] (1.5,0.3) -- (3,1.5);
  \draw[-{Stealth},thick,purple] (-1.5,-0.3) -- (-3,-1.5);
  \draw[-{Stealth},thick,orange] (1.5,-0.3) -- (3,-1.5);
\end{tikzpicture}
\end{center}

\subsection{Practical Strategies for NP-Hard Problems}

When you encounter an NP-complete problem in practice:

\begin{center}
\begin{tabular}{>{\bfseries}p{3.5cm} p{5cm} p{4cm}}
\toprule
Strategy & Idea & Example \\
\midrule
Approximation Algorithms & Find near-optimal solutions in poly time & $2$-approx for Vertex Cover \\
Heuristics & Good practical performance, no guarantee & Simulated annealing \\
Parameterized Algorithms & Poly in $n$, exponential only in a small parameter & FPT for treewidth \\
Special Cases & Exploit structure for poly algorithms & 2-SAT, planar graphs \\
Randomized Algorithms & Probabilistic guarantees & Random restart SAT \\
\bottomrule
\end{tabular}
\end{center}

\subsection{Real-World Applications of SAT}

Modern SAT solvers (DPLL, CDCL) can handle industrial instances with millions of variables:

\begin{itemize}[leftmargin=2em]
  \item \textbf{Hardware verification}: Intel, AMD use SAT to verify chip designs.
  \item \textbf{Software model checking}: Checking program correctness.
  \item \textbf{Cryptanalysis}: Finding weaknesses in cryptographic protocols.
  \item \textbf{Planning}: AI planning in robotics and scheduling.
  \item \textbf{Bioinformatics}: Haplotype phasing, protein folding.
\end{itemize}

\newpage

% ══════════════════════════════════════════════
\section{Summary and Key Takeaways}
% ══════════════════════════════════════════════

\begin{tcolorbox}[
  enhanced, breakable,
  colback=lightblue, colframe=darkblue, fonttitle=\bfseries\large,
  title={The Cook--Levin Theorem: Complete Summary},
  attach boxed title to top left={yshift=-3mm,xshift=4mm},
  boxed title style={colback=darkblue}
]
\begin{enumerate}[leftmargin=1.5em, itemsep=0.8em]
  \item \textbf{P vs NP}: The central question --- can every efficiently-verifiable problem 
    be efficiently solved?

  \item \textbf{NP Definition}: A problem is in NP if ``yes'' instances have short, 
    efficiently-verifiable certificates.

  \item \textbf{Polynomial Reduction} ($A \leqp B$): A poly-time translation preserving 
    yes/no answers. Used to compare difficulty.

  \item \textbf{NP-Completeness}: A problem is NP-complete if it is in NP \emph{and} 
    every NP problem reduces to it.

  \item \textbf{SAT is NP-complete} (Cook--Levin):
    \begin{itemize}
      \item[$\bullet$] SAT $\in$ NP: assignments are poly-size certificates.
      \item[$\bullet$] Every NP problem reduces to SAT: encode the TM's computation 
        as a CNF formula via the tableau construction.
    \end{itemize}

  \item \textbf{Proof key steps}: 
    NP problem $\Rightarrow$ poly-time verifier $\Rightarrow$ TM tableau 
    $\Rightarrow$ four groups of clauses $\Rightarrow$ poly-size CNF formula.

  \item \textbf{Implication}: If SAT $\in$ P, then P $=$ NP.
    If SAT $\notin$ P, then P $\neq$ NP.

  \item \textbf{Cascade}: Once SAT is NP-complete, we can prove hundreds of 
    other problems NP-complete via chain reductions.
\end{enumerate}
\end{tcolorbox}

\bigskip

\begin{figure}[H]
\centering
\begin{tikzpicture}[scale=0.85, every node/.style={font=\small}]
  % Main concept map
  \node[draw=darkblue,fill=lightblue,rounded corners,minimum width=2.5cm,
        minimum height=0.8cm,align=center,font=\bfseries] (TM) at (0,5) 
    {Turing Machine};
  \node[draw=darkgreen,fill=lightgreen,rounded corners,minimum width=2.5cm,
        minimum height=0.8cm,align=center,font=\bfseries] (NP) at (0,3.2) 
    {\NP\ class};
  \node[draw=purple,fill=purple!15,rounded corners,minimum width=2.5cm,
        minimum height=0.8cm,align=center,font=\bfseries] (red) at (-3.5,1.5) 
    {Poly-Time Reduction};
  \node[draw=orange,fill=orange!15,rounded corners,minimum width=2.5cm,
        minimum height=0.8cm,align=center,font=\bfseries] (tab) at (3.5,1.5) 
    {TM Tableau};
  \node[draw=darkred,fill=lightred,rounded corners,minimum width=2.5cm,
        minimum height=0.8cm,align=center,font=\bfseries] (SAT) at (0,0) 
    {\SAT\ (NP-complete)};
  \node[draw=teal,fill=teal!15,rounded corners,minimum width=2.5cm,
        minimum height=0.8cm,align=center,font=\bfseries] (chain) at (0,-1.8)
    {Chain Reductions};
  \node[draw=darkgray,fill=lightgray,rounded corners,minimum width=2.5cm,
        minimum height=0.8cm,align=center,font=\bfseries] (nph) at (0,-3.5)
    {NP-Hard Problems};

  \draw[-{Stealth},thick,darkblue] (TM) -- (NP) node[midway,right] {defines};
  \draw[-{Stealth},thick,darkgreen] (NP) -- (red) node[midway,left,sloped] {uses};
  \draw[-{Stealth},thick,darkgreen] (NP) -- (tab) node[midway,right,sloped] {encodes to};
  \draw[-{Stealth},thick,purple] (red) -- (SAT) node[midway,left,sloped] {proves};
  \draw[-{Stealth},thick,orange] (tab) -- (SAT) node[midway,right,sloped] {gives};
  \draw[-{Stealth},thick,darkred] (SAT) -- (chain) node[midway,right] {enables};
  \draw[-{Stealth},thick,teal] (chain) -- (nph) node[midway,right] {classifies};
\end{tikzpicture}
\caption{Concept map of the Cook--Levin theorem and its connections.}
\end{figure}

\newpage

% ══════════════════════════════════════════════
\section{Practice Questions}
% ══════════════════════════════════════════════

\subsection{Conceptual Questions}

\begin{exercise}{P vs NP Understanding}{pnp1}
\begin{enumerate}[label=(\alph*),leftmargin=2em]
  \item Explain in your own words why NP is defined in terms of \emph{verification}, not computation.
  \item Give an example of a problem that is:
    \begin{enumerate}[label=(\roman*)]
      \item In P (and hence in NP), and
      \item In NP but (as far as we know) not in P.
    \end{enumerate}
  \item If someone proves that $\P = \NP$, what would be the consequence for cryptography (e.g., RSA encryption)?
\end{enumerate}
\end{exercise}

\begin{exercise}{Reduction Mechanics}{red1}
Let $A \leqp B$ and $B \leqp A$.
\begin{enumerate}[label=(\alph*),leftmargin=2em]
  \item Does this imply $A = B$? Explain.
  \item Does this imply $A$ and $B$ have the same complexity? Explain.
  \item If $A \in \P$, does it follow that $B \in \P$? Why or why not?
\end{enumerate}
\end{exercise}

\begin{exercise}{NP-Completeness Definition}{npc1}
For each statement, decide if it is \textbf{True} or \textbf{False} and justify your answer.
\begin{enumerate}[label=(\alph*),leftmargin=2em]
  \item Every NP-complete problem is in NP.
  \item If problem $A$ is NP-complete and $A \leqp B$, then $B$ is NP-hard.
  \item If $\P \neq \NP$, then no NP-complete problem is in P.
  \item SAT is the only NP-complete problem.
  \item Showing a problem is NP-hard is sufficient to show it is NP-complete.
\end{enumerate}
\end{exercise}

\subsection{SAT and Formula Questions}

\begin{exercise}{SAT Evaluation}{sat_eval}
For each formula, determine if it is satisfiable, and if so, give a satisfying assignment:
\begin{enumerate}[label=(\alph*),leftmargin=2em]
  \item $\phi_1 = (x \lor y) \land (\neg x \lor y) \land (\neg y)$
  \item $\phi_2 = (x \lor y \lor z) \land (\neg x \lor \neg y) \land (\neg y \lor \neg z) \land (\neg x \lor \neg z)$
  \item $\phi_3 = (x_1 \lor \neg x_2) \land (\neg x_1 \lor x_2) \land (x_1 \lor x_2) \land (\neg x_1 \lor \neg x_2)$
\end{enumerate}
\end{exercise}

\begin{exercise}{CNF Conversion}{cnf1}
Convert the following formula to CNF (and then check if it is satisfiable):
\[
  \phi = (x_1 \Rightarrow x_2) \land (x_2 \Rightarrow x_3) \land \neg x_3
\]
\textit{Hint: $a \Rightarrow b \equiv \neg a \lor b$.}
\end{exercise}

\begin{exercise}{Long Clause to 3-SAT}{3sat1}
Convert the clause $(x_1 \lor x_2 \lor x_3 \lor x_4 \lor x_5 \lor x_6)$ to an 
equivalent 3-CNF formula using the method from Section 6.2.
How many new variables and clauses do you introduce?
\end{exercise}

\subsection{Proof-Oriented Questions}

\begin{exercise}{Tableau Construction}{tableau1}
Consider the TM $M$ with states $\{q_0, q_1, q_\text{acc}\}$, alphabet $\{0,1,\sqcup\}$,
and transitions:
\[
  \delta(q_0, 0) = (q_1, 1, R),\quad \delta(q_1, \sqcup) = (q_\text{acc}, \sqcup, R)
\]
Input: $x = 0$. Take time bound $T = 3$, tape width $W = 3$.
\begin{enumerate}[label=(\alph*),leftmargin=2em]
  \item Draw the complete $4 \times 3$ tableau.
  \item Write out all variables $x_{i,j,s}$ that are set to 1.
  \item Write the $\phi_\text{start}$ clauses for this instance.
  \item Write the $\phi_\text{accept}$ clause.
\end{enumerate}
\end{exercise}

\begin{exercise}{Cook--Levin Key Steps}{cl_steps}
Answer the following about the Cook--Levin proof:
\begin{enumerate}[label=(\alph*),leftmargin=2em]
  \item What is the purpose of the four groups of clauses 
    ($\phi_\text{cell}$, $\phi_\text{start}$, $\phi_\text{move}$, $\phi_\text{accept}$)?
  \item Why is the tableau size $q(n) \times q(n)$ (not $q(n) \times n$)?
  \item What does the satisfying assignment of $\phi_x$ ``encode''?
  \item If the input $x \notin L$, why is $\phi_x$ unsatisfiable?
  \item Why is the reduction computable in polynomial time?
\end{enumerate}
\end{exercise}

\begin{exercise}{Show NP-Completeness}{show_npc}
Suppose you are given a new problem $\Pi$. To show $\Pi$ is NP-complete, 
what two things must you prove? For each, describe the general approach:
\begin{enumerate}[label=(\alph*),leftmargin=2em]
  \item Show $\Pi \in \NP$.
  \item Show $\Pi$ is NP-hard (i.e., show $\SAT \leqp \Pi$ or some known NP-complete problem $\leqp \Pi$).
\end{enumerate}
\end{exercise}

\subsection{Challenging Problems}

\begin{exercise}{2-SAT is in P}{twosat}
\textbf{2-SAT} is the restriction of SAT to clauses with exactly 2 literals.
\begin{enumerate}[label=(\alph*),leftmargin=2em]
  \item Show 2-SAT $\in$ NP.
  \item Is 2-SAT NP-complete? (No, it is actually in P.) 
    Describe intuitively why having only 2 literals per clause makes the problem easier.
  \item The polynomial-time algorithm for 2-SAT uses \textbf{strongly connected components}. 
    Look this up and write a 3--5 sentence explanation.
\end{enumerate}
\end{exercise}

\begin{exercise}{Clique is NP-Complete}{clique_npc}
Show that \CLIQUE\ is NP-complete by proving $\TSAT \leqp \CLIQUE$.

\medskip
\textit{Hint:} Given a 3-CNF formula $\phi$ with $k$ clauses, construct a graph $G$ 
with $3k$ vertices (one per literal per clause). Add an edge between $u$ and $v$ if 
they are in different clauses and are not contradictory ($u \neq \neg v$).

Show: $\phi$ is satisfiable $\iff$ $G$ has a clique of size $k$.
\end{exercise}

\begin{exercise}{Reflection and Application}{reflect}
\begin{enumerate}[label=(\alph*),leftmargin=2em]
  \item A company tells you ``we have a polynomial-time algorithm for solving any 
    Boolean formula in 3-CNF form.'' What would this imply about P vs NP?
  \item You are designing a scheduling system for 1000 employees and discover 
    your problem is NP-complete. What practical strategies would you recommend? 
    List at least three.
  \item The Clay Mathematics Institute offers \$1,000,000 for solving P vs NP. 
    Why do most researchers believe $\P \neq \NP$? What evidence supports this?
\end{enumerate}
\end{exercise}

\newpage

% ══════════════════════════════════════════════
\section{Selected Answers and Hints}
% ══════════════════════════════════════════════

\subsection*{Exercise 7.2.1 (SAT Evaluation)}

\noindent\textbf{(a)} $\phi_1 = (x \lor y) \land (\neg x \lor y) \land (\neg y)$
\begin{itemize}[leftmargin=2em]
  \item Clause 3 forces $y = 0$.
  \item Then clause 1: $(x \lor 0) = x$ forces $x = 1$.
  \item Check clause 2: $(\neg 1 \lor 0) = (0 \lor 0) = 0$. \textbf{UNSAT.}
\end{itemize}

\noindent\textbf{(b)} $\phi_2$: try $x = 1, y = 0, z = 0$:
Clause 1: $(1 \lor 0 \lor 0)=1$. Clause 2: $(0 \lor 1)=1$. Clause 3: $(1 \lor 1)=1$. 
Clause 4: $(0 \lor 1)=1$. \textbf{SAT.}

\subsection*{Exercise 7.2.2 (CNF Conversion)}
$\phi = (\neg x_1 \lor x_2) \land (\neg x_2 \lor x_3) \land \neg x_3$

Clause 3 forces $x_3 = 0$. Clause 2 forces $x_2 = 0$ (since $\neg 0 = 1$, actually: 
$(\neg x_2 \lor 0)$ forces $x_2 = 0$). Clause 1: $(\neg x_1 \lor 0)$ forces $x_1 = 0$. 
Assignment $x_1=0, x_2=0, x_3=0$ satisfies all clauses. \textbf{SAT.}

\subsection*{Exercise 7.3.1 (Tableau Construction)}
The 4$\times$3 tableau:
\begin{center}
\begin{tabular}{c|ccc}
\toprule
$t$ & Cell 0 & Cell 1 & Cell 2 \\
\midrule
0 & $(q_0,0)$ & $\sqcup$ & $\sqcup$ \\
1 & $1$ & $(q_1,\sqcup)$ & $\sqcup$ \\
2 & $1$ & $\sqcup$ & $(q_\text{acc},\sqcup)$ \\
3 & $1$ & $\sqcup$ & $(q_\text{acc},\sqcup)$ \\
\bottomrule
\end{tabular}
\end{center}

\subsection*{Hints for Challenging Problems}

\textbf{Exercise 7.4.2 (Clique):}
\begin{itemize}[leftmargin=2em]
  \item \textit{NP part:} A clique of $k$ vertices is a certificate of size $O(k)$.
  \item \textit{Forward direction:} If $\phi$ is satisfied by assignment $\alpha$, each clause 
    has at least one true literal. Pick one true literal per clause. These $k$ literal-nodes 
    form a clique (they're in different clauses, and non-contradictory since $\alpha$ is consistent).
  \item \textit{Backward direction:} A $k$-clique has one node per clause (pigeonhole: $k$ clauses, 
    $k$ nodes, all different clauses). Setting each chosen literal to true gives a consistent 
    satisfying assignment.
\end{itemize}

\newpage

% ══════════════════════════════════════════════
\appendix
\section{Quick Reference Card}
% ══════════════════════════════════════════════

\begin{center}
\renewcommand{\arraystretch}{1.5}
\begin{tabular}{>{\bfseries\color{darkblue}}p{3.5cm} p{8cm}}
\toprule
\textcolor{darkblue}{Concept} & \textcolor{darkblue}{Definition / Key Fact} \\
\midrule
$\P$ & Poly-time decidable by a deterministic TM \\
$\NP$ & Poly-time verifiable; equivalently, decidable by a nondeterministic TM in poly time \\
$\P \subseteq \NP$ & Every poly-time algorithm is also a poly-time verifier \\
$A \leqp B$ & $A$ poly-time reduces to $B$; if $B \in \P$ then $A \in \P$ \\
NP-hard & Every NP problem reduces to it \\
NP-complete & NP-hard AND in NP \\
Cook--Levin & SAT is NP-complete \\
3-SAT & SAT restricted to 3 literals/clause; also NP-complete \\
Tableau & $q(n) \times q(n)$ grid encoding a TM computation \\
$\phi_\text{cell}$ & Each tableau cell has exactly one symbol \\
$\phi_\text{start}$ & First row = initial configuration \\
$\phi_\text{move}$ & Each row = legal TM transition from previous \\
$\phi_\text{accept}$ & Final row contains accept state \\
Certificate size & At most polynomial in input length \\
If $\P = \NP$ & All NP-complete problems are in P \\
\bottomrule
\end{tabular}
\end{center}

\vspace{1cm}

\begin{center}
\begin{tikzpicture}[scale=0.9]
  \fill[gray!15,rounded corners=20pt] (-5.2,-3.5) rectangle (5.2,3.5);
  \node[font=\small\bfseries,color=darkgray] at (0,3.1) {EXPTIME};
  
  \fill[midblue!30,rounded corners=14pt] (-4.5,-2.8) rectangle (4.5,2.8);
  \node[font=\bfseries,color=darkblue] at (0,2.4) {\textbf{NP}};
  
  % NP-complete region
  \fill[lightred,rounded corners=8pt] (-4.3,-2.6) rectangle (-0.5,2.6);
  \node[font=\small\bfseries,color=darkred,align=center] at (-2.4,0) 
    {NP-complete\\SAT, 3-SAT,\\Clique, Ham-Path,\\\ldots};
  
  \fill[lightblue,rounded corners=10pt] (0,-2.4) rectangle (4.3,2.4);
  \node[font=\bfseries,color=darkblue] at (2.2,1.8) {\P};
  \node[font=\small,color=darkblue,align=center] at (2.2,0)
    {Sorting, Shortest\\Path, 2-SAT,\\Primality,\ldots};
  
  \draw[darkred,very thick,dashed] (-0.5,-2.5) -- (-0.5,2.5);
  \node[font=\tiny\bfseries,color=darkred,rotate=90] at (-0.7,0) {P=NP boundary?};
\end{tikzpicture}
\end{center}

\vspace{1cm}

\begin{tcolorbox}[colback=lightyellow,colframe=orange,fonttitle=\bfseries,
  title={The One-Sentence Takeaway}]
\centering\large
\textit{The Cook--Levin theorem shows that SAT is the ``hardest'' problem in NP ---\\
solving it in polynomial time would solve \textbf{all} of NP in polynomial time.}
\end{tcolorbox}

\end{document}
